\documentclass{article}
\usepackage{amsmath, amssymb, amsfonts}

\newcommand{\op}{\prime}
\title{Ionic Crystals}
\begin{document}
\author{Amey Joshi}
\maketitle
Ionic crystals, like the crystals of inert elements, are formed by atoms with
completely filled outermost orbitals. However, in the case of ionic crystals, 
the constituents are charged and the dominant force between then is 
electrostatic rather than Van der Waals. The force between charges is usually
much stronger than that between dipoles. Quite like the case of inert element
crystals, there is very little overlap between the electronic clouds of the 
constituents. This fact is well supported by the electronic density maps of
the crystals. The electrostatic energy binding the crystal is called Madelung
energy after its inventor.

If $U_{ij}$ is the energy of interaction between ions $i$ and $j$ then 
\begin{equation}\label{e1}
U_{ij} = \lambda e^{-r_{ij}/\rho} \pm \frac{1}{4\pi\epsilon_0}\frac{q^2}{r_{ij}},
\end{equation}
where the second term is the Coulomb potential. The first term represents 
a repulsive potential with experimentally adjustable parameters $\lambda$ and
$\rho$. This term prevents $U_{ij}$ from getting minimised by letting $r_{ij}$
be zero. It is a simple form of Pauli exclusion principle that stops the fully
filled outer orbitals of the ions from overlapping each other although the ions
have a Coulomb attraction between them.

The total energy of the crystal is
\begin{equation}\label{e2}
U_\text{tot} = \frac{1}{2}\sideset{}{^\prime}\sum_{ij}U_{ij},
\end{equation}
where the prime on the sum indicates that we are dropping the $i = j$ term and
the factor of $1/2$ corrects the double counting of every pair of ions.

We now assume that the total energy of an ionic crystal is $NU_i$, where $N$ is
the number of atoms in it and $U_i$ is the potential energy of a single ion in
it. This approximation is valid for large enough crystals
because the number of atoms on the surface is much smaller than those in the 
bulk. The expression for $U_i$ is quite different for the ions at the surface 
than those in the interior. To illustrate this point we will consider a crystal
of metallic sodium. We will show that it is \emph{not} an ionic crystal. However
we use it as an example only to compare the number of atoms on the surface with
those in the bulk. Consider a cubical crystal of side $1$ mm. Its volume is
$1$ mm$^3$ or $10^{-3}$ cm$^3$. The density of Na is $0.97$ g/cm$^3$. Therefore,
its mass is $M = 0.97 \times 10^{-3}$. It atomic weight is $23$ g. Therefore, a
cube of mass $M$ has
\[
N = \frac{M \times N_A}{23} \approx 2.54 \times 10^{19}
\]
atoms. An atom of Na has a diameter of roughly $4 \times 10^{-8}$ cm. Therefore,
the number of atoms on each face is
\[
N_s = \frac{10^{-2}}{4 \times 10^{-8}} = 0.25 \times 10^6.
\]
Therefore, the ratio of atoms on all six surfaces to the total number of atoms
is
\[
\frac{1.5 \times 10^6}{2.54 \times 10^{19}} = O(10^{-12}),
\]
a very small fraction. Since the repulsive term in \eqref{e2} decays rapidly
with distance, it is relevant only for the nearest neighbours. If there are $z$
nearest neighbours and $R$ is the distance between them then we can write
\begin{equation}\label{e3}
U_i = z\lambda e^{-R/\rho} - z\frac{1}{4\pi\epsilon_0}\frac{q^2}{R}  \pm 
\frac{1}{4\pi\epsilon_0}\frac{q^2}{R} 
\sideset{}{^\prime}\sum_{ij}\frac{1}{p_{ij}},
\end{equation}
where $p_{ij} = r_{ij}/R$. We write it as
\begin{equation}\label{e4}
U_i = z\lambda e^{-R/\rho} - \frac{1}{4\pi\epsilon_0}\frac{q^2}{R}\alpha,
\end{equation}
where
\begin{equation}\label{e5}
\alpha = z \mp \sideset{}{^\prime}\sum_{ij}\frac{1}{p_{ij}},
\end{equation}
is called the Madelung constant for the crystal. It is a geometrical metric
associated with the crystal. We will take it as a given quantity. The total
energy of the crystal is
\begin{equation}\label{e6}
U_\text{tot} = \frac{N}{2}U_i = \frac{N}{2}
\left(z\lambda e^{-R/\rho} -\frac{1}{4\pi\epsilon_0}\frac{q^2}{R}\alpha\right).
\end{equation}
The crystal settles with a nearest neighbour distance that minimises its energy.
Now,
\[
\frac{dU_\text{tot}}{dR} = \frac{N}{2}\left(-\frac{z\lambda}{\rho}e^{-R/\rho} + 
\frac{1}{4\pi\epsilon_0} \frac{q^2}{R^2}\alpha\right)
\]
so that at extremum, we have
\begin{equation}\label{e7}
z\lambda  e^{-R_0/\rho} = \frac{1}{4\pi\epsilon_0} \frac{q^2\alpha\rho}{R_0^2}.
\end{equation}
Substituting this in \eqref{e6} we get the equilibrium energy of the crystal
\begin{equation}\label{e8}
U_\text{tot} = -\frac{N}{2}\frac{q^2}{4\pi\epsilon_0}\left[\frac{\alpha}{R_0}
\left(1 - \frac{\rho}{R_0}\right)\right].
\end{equation}
The terms in the brackets specific to the crystal structure. If we write this
equation as 
\begin{equation}\label{e9}
\frac{U_\text{tot}}{N} = -\frac{1}{2}\frac{q^2}{4\pi\epsilon_0}
\left[\frac{\alpha}{R_0} \left(1 - \frac{\rho}{R_0}\right)\right].
\end{equation}
so that it gives an expression for the electrostatic energy per ion, then the
expression on the right hand side follows only from the crystal's geometry.
The parameter $\rho$ is usually of the order of $0.1 R_0$ $U_\text{tot}$ is
guaranteed to be negative. The smallness of $\rho$ also guarantees that the 
exponential term in \eqref{e2} is indeed quite small. At $r_{ij} = \rho$, it
decays to $1/e$ times that its value at $r_{ij} = 0$. As $\rho$ is expected to
be a small number, a rough estimate of electrostatic energy of an ionic crystal
is
\begin{equation}\label{e10}
U_\text{tot} = -\frac{N}{2}\frac{q^2}{4\pi\epsilon_0}\frac{\alpha}{R}.
\end{equation}
It is also called Madelung energy.

\section{Problems}
\begin{enumerate}
\item Ionic crystal of metallic sodium: Let us examine the feasibility of 
metallic sodium forming an ionic crystal. We need two reactions for that to
happen:
\begin{eqnarray*}
\text{Na} + 5.14 \text{eV} &\rightarrow& \text{Na}^+ + e^- \\
\text{Na} + e^- &\rightarrow& \text{Na}^- + 0.78 \text{eV}
\end{eqnarray*}
In the first reaction, neutral Na atom ionises after acquiring $5.14$ eV while
in the second one a neutral Na atom release $0.78$ eV after accepting an 
electron. The energy in the second equation is called the electron affinity of
the atom. The net of these reactions is
\[
2\text{Na} + 4.36 \text{eV} \rightarrow \text{Na}^+ + \text{Na}^- 
\]
Equivalently, the energy required \emph{per atom} to make this happen is $4.36
/2 = 2.18$ eV.

If we assume that it crystallises like NaCl then the Madelung constant is
$\alpha = 1.747565$. If we assume that $\rho \approx 0.1$ then the energy per
ion, using \eqref{e9}, is
\[
\frac{U_\text{tot}}{N} = -\frac{1}{2} \times (1.6 \times 10^{-19}) \times 
(9 \times 10^9) \times \left[0.9 \times \frac{1.747565}{3.66 \times 10^{-10}}
\right] \approx -3.095 \text{eV},
\]
where we have used the fact that in an fcc crystal, $R_0 = 0.866 a$, $a$ being
the lattice constant. It is $4.225$ \AA for NaCl.

Therefore, the net drop in energy by forming an ionic crystal per atom is
$2.18 - 3.095 \approx -0.915$ eV. The cohesive energy of a crystal of Na metal
is $1.113$ eV per atom. That is, the potential energy of an atom in Na crystal
is $-1.113$ eV. This is less than the energy of the hypothetical Na$^+$Na$^-$
crystal. Thus, on the basis of energy arguments alone we can conclude that the
ionic crystal structure is not a favourable arrangement.

\item Linear ionic crystal: Consider a line of ions of alternating charges
separated by a distance $R$ and let the repulsive part of the interaction by
$A/R^n$ where $A > 0$ is a constant and $n$ is a positive integer. The potential
energy between ions $i$ and $j$ is
\[
U_{ij}=\frac{A}{R^n} - \frac{1}{4\pi\epsilon_0}\sum_{k=-\infty, k \ne 0}^\infty 
\frac{q^2}{|kR|} = \frac{A}{R^n} - \frac{2\ln 2}{4\pi\epsilon_0}\frac{q^2}{R}
\]
If there are $N$ ions in the crystal, then ignoring the surface effects, the
energy of the crystal is
\[
U = \frac{1}{2}\sideset{}{^\prime}\sum_{ij} U_{ij} = \frac{N}{2}\left(\frac{A}
{R^n} - \frac{2\ln 2}{4\pi\epsilon_0}\frac{q^2}{R}\right).
\]
Differentiating with respect to $R$,
\[
\frac{dU}{dR} = \frac{N}{2}\left(-n\frac{A}{R^{n+1}} + 
\frac{2\ln 2}{4\pi\epsilon_0}\frac{q^2}{R^2} \right)
\]
so that the point of extremum is
\[
R_0^{n-1} = \frac{nA}{(2\ln 2) q^2} \Rightarrow \frac{A}{R_0^{n}} = 
\frac{(2\ln 2)q^2}{nR_0}.
\]
At the extremum,
\[
U(R_0) = -N\frac{q^2\ln 2}{4\pi\epsilon_0R_0}\left(1 - \frac{1}{n}\right).
\]
This expression differs from the one in Kittel because we have a factor of 
$1/2$ to take care of double counting when there are $N$ atoms and not $N$ 
pairs as in the book.

A Taylor expansion of $U$ is
\[
U(R + h) = U(R) + hU^\prime(R) + \frac{h^2}{2}U^{\prime\prime}(R) + O(h^3).
\]
At $R = R_0$ the second term vanishes and we have
\[
U(R_0 + h) = U(R_0) + \frac{h^2}{2}U^{\prime\prime}(R_0).
\]
From the expression for $dU/dR$, we have
\[
U^{\prime\prime}(R) = \frac{N}{2}\left(n(n+1)\frac{A}{R^{n+2}} - 
4\ln 2\frac{q^2}{R^3}\right)
\]
When $R = R_0$, 
\[
\frac{nA}{R_0^{n+2}} = \frac{2\ln 2}{4\pi\epsilon_0}\frac{q^2}{R_0^3},
\]
so that
\[
U^{\prime\prime}(R_0) = \frac{N}{4\pi\epsilon_0}\frac{q^2\ln 2}{R_0^3}(n - 1).
\]
For a crystal of unit length, $NR_0 = 1$ and the above expression becomes
\[
U^{\prime\prime}(R_0) = \frac{1}{4\pi\epsilon_0}\frac{q^2\ln 2}{R_0^4}(n - 1).
\]
The force acting on the crystal is
\[
F = -\frac{U(R_0 + h) - U(R_0)}{h}
\]
and the work done in compressing it is $-Fh$ that is $U(R_0+h)-U(R_0)$, that is,
\[
\frac{1}{2}\frac{1}{4\pi\epsilon_0}\frac{q^2\ln 2}{R_0^4}(n - 1)h^2.
\]
The change in interatomic distance is expressed as $R_0 = R_0(1 - \delta)$ so 
that $h = R_0\delta$. The above expression then reduces to
\[
\frac{1}{2}\frac{1}{4\pi\epsilon_0}\frac{q^2\ln 2}{R_0^2}(n - 1)\delta^2.
\]
I think there is a typo in the book. The quantity work per unit length must have
dimensions of force.
\end{enumerate}
\end{document}
